%%%%%%%%%%%%%%%%%%%%%%%%%%%%%%%%%%%%%%%%%%%%%%%%%%%%%%%%%%%%%%%%%%%%%%%%%%%%%%%%
% Obxectivo: Lista de termos empregados no documento,                          %
%            xunto cos seus respectivos significados.                          %
%%%%%%%%%%%%%%%%%%%%%%%%%%%%%%%%%%%%%%%%%%%%%%%%%%%%%%%%%%%%%%%%%%%%%%%%%%%%%%%%

\newglossaryentry{bytecode}{
  name=bytecode,
  description={Código independente da máquina que xeran compiladores de determinadas linguaxes (Java, Erlang,\dots) e que é executado polo correspondente intérprete.}
}
\newglossaryentry{Spotify}{
	name=Spotify,
	description={Servicio dixital de reprodución de música, podcast e vídeos.}
}
\newglossaryentry{frontend}{
	name=frontend,
	description={O frontend é a parte da aplicación que interactúa cos usuarios, é coñecida coma o lado do cliente.}
}
\newglossaryentry{backend}{
	name=backend,
	description={O backend é a parte do desenvolvemento web que se encaga da lóxica para que as aplicacións funcionen.}
}
\newglossaryentry{ionic}{
	name=ionic,
	description={Conxunto de ferramentas para desenvolver aplicacións de alta calidade e multiplataforma.}
}
\newglossaryentry{django}{
	name=django,
	description={Framework web de alto nivel en Python que facilita o desenvolvemento rápido e limpio.}
}
\newglossaryentry{framework}{
	name=framework,
	description={Entorno de traballo con un conxunto de conceptos, prácticas e criterios estandarizado para resolver un tipo de problemática particular.}
}
\newglossaryentry{API}{
		name=API,
		description={Conxunto de definicións e protocolos que se utiliza para desenvolver e integrar software das aplicacións.}
}
\newglossaryentry{TypeScript}{
	name=TypeScript,
	description={Linguaxe de programación libre e de código aberto mantido por Microsoft.}
}
\newglossaryentry{Scrum}{
	name=Scrum,
	description={É un proceso no que se aplica de maneira regular un conxunto de boas prácticas para traballar colaborativamente e obter o mellor resultado posible dun proxecto.}
}
